\begin{landscape}
\begin{table}[p]
\centering
\small
\begin{threeparttable}
\caption{RQ2: accessibility before and after area socioeconomic context}\label{tab:rq2-area-context}
\begin{tabular}{>{\raggedright\arraybackslash}p{2.0cm}>{\raggedright\arraybackslash}p{6.0cm}rrrrrr}
\toprule
Estimator & Specification & $N$ & M0 RMSE & ME RMSE & MW RMSE & ME gain & MW gain \\
\midrule
OLS & Base specification & 12,412 & 0.33446 & 0.32776 & 0.33053 & 0.00670 & 0.00393 \\
OLS & Area-context-adjusted core specification & 12,412 & 0.32407 & 0.32498 & 0.32592 & -0.00091 & -0.00185 \\
XGBoost & Base specification & 12,412 & 0.33277 & 0.32795 & 0.32976 & 0.00482 & 0.00301 \\
XGBoost & Area-context-adjusted core specification & 12,412 & 0.32675 & 0.32670 & 0.32731 & 0.00004 & -0.00056 \\
\bottomrule
\end{tabular}
\begin{tablenotes}[flushleft]\footnotesize
\item Pooled leave-one-area-out log-price RMSE on the same listings and saved folds. Base = Phase 9 listing, host/rental, municipality and City Hall centrality controls. Area-context-adjusted core = base plus log gross-polygon population density and disposable income; this planned contextual specification was implemented after Phase 9, not frozen at Phase 7. Positive gain means lower error than the matching M0. Mixed Copenhagen district/Frederiksberg municipality context, income definitions and polygon areas remain provisional.
\end{tablenotes}
\end{threeparttable}
\end{table}
\end{landscape}
